\documentclass[a4paper,11pt]{article}
\usepackage[utf8]{inputenc}
\usepackage[T1]{fontenc}
\usepackage[french]{babel}
\usepackage[left=2cm,right=2cm,top=2cm,bottom=2cm]{geometry}
\usepackage{xcolor}
\usepackage{hyperref}
\usepackage{fontawesome5}
\usepackage{parskip}

% --- Couleurs Enedis (Bleu énergie) ---
\definecolor{primary}{RGB}{0, 87, 184}
\hypersetup{colorlinks=true, urlcolor=primary}

\begin{document}

% --- EXPÉDITEUR ---
\noindent
\begin{minipage}[t]{0.5\textwidth}
    \textbf{Loïc Aron MBASSI EWOLO}\\
    Élève-Ingénieur Bac+5 -- Double diplôme ENSEA\\[3pt]
    \faMapMarker\ Cergy, France\\
    \faPhone\ (+33) 7 59 52 33 98\\
    \faEnvelope\ \href{mailto:loicaron.mbassiewolo@ensea.fr}{loicaron.mbassiewolo@ensea.fr}
\end{minipage}
\hfill
% --- DATE ---
\begin{minipage}[t]{0.4\textwidth}
    \raggedleft
    Cergy, le 2 mars 2026
\end{minipage}

\vspace{1.2cm}

% --- DESTINATAIRE ---
\hfill
\begin{minipage}[t]{0.5\textwidth}
    \raggedleft
    \textbf{Enedis}\\
    Service Recrutement\\
    Nice, France
\end{minipage}

\vspace{1cm}

\noindent\textbf{Objet :} Candidature en alternance -- Concepteur Développeur F/H

\vspace{0.8cm}

Madame, Monsieur,

Trois ans de développement PHP en production, six plateformes web livrées en autonomie, et une conviction : la qualité du code commence par une analyse rigoureuse du besoin métier. C'est avec cet état d'esprit que je souhaite rejoindre Enedis en alternance en tant que Concepteur Développeur, pour contribuer à la conception et à l'évolution de vos applications web au service du réseau électrique.

En double diplôme à l'ENSEA (Systèmes Embarqués, Signal, IA) après un cursus d'ingénieur en Mathématiques Appliquées et Génie Logiciel à l'École Polytechnique de Yaoundé, je combine une solide formation académique avec une expérience professionnelle concrète en PHP/Laravel. Chez SOSDOCTA, j'ai développé et maintenu pendant deux ans une plateforme de télémédecine complète : conception du backend, intégration frontend, système de réservation, gestion des données médicales conformes au RGPD. Cette expérience m'a appris à rédiger des spécifications techniques claires, à structurer les données pour répondre à des besoins métier complexes, et à assurer la maintenance évolutive d'une application en production.

Plus récemment, chez CSC dans le secteur fintech, j'ai conçu le backend d'une plateforme de transactions financières sécurisée avec Laravel, MySQL, Docker et Redis. La rédaction de cahiers des charges pour les intégrations tierces et la gestion de transactions concurrentes m'ont confirmé l'importance d'une documentation technique rigoureuse. En parallèle, j'ai développé un outil open source (Laravel API Generator, +30 étoiles GitHub) qui génère automatiquement une architecture backend complète depuis des diagrammes UML/XML, combinant analyse de données structurées et programmation. Mon intérêt pour le traitement des données s'est aussi concrétisé lors d'un hackathon national (1\textsuperscript{er} prix CITS), où j'ai modélisé et analysé des données IoT pour optimiser la collecte de déchets urbains.

La mission d'Enedis, acteur central de la transition énergétique, donne un sens concret au développement logiciel. Contribuer à des applications qui soutiennent la gestion du réseau électrique français correspond à ma volonté de mettre mes compétences techniques au service d'un enjeu collectif. Je suis disponible dès septembre 2026 pour une alternance et mobile pour rejoindre le site de Nice.

Je reste à votre disposition pour un entretien afin de discuter de ma candidature.

Je vous prie d'agréer, Madame, Monsieur, l'expression de mes salutations distinguées.

\vspace{0.8cm}

\textbf{Loïc Aron MBASSI EWOLO}\\
\textit{Élève-Ingénieur ENSEA / Polytechnique Yaoundé}

\end{document}
